\documentclass[11pt,a4paper]{article}
\usepackage[T1]{fontenc}
\usepackage[utf8]{inputenc}
\usepackage[english]{babel}
\usepackage{amsmath,amssymb}
\usepackage{geometry}
\geometry{margin=2.3cm}
\usepackage{booktabs}
\usepackage{graphicx}
\usepackage{lineno}
\usepackage{hyperref}
\hypersetup{colorlinks=true,linkcolor=blue,citecolor=blue,urlcolor=blue}

\title{Emergent Time from Correlations:\\
A Unified Framework for Cosmological Evolution}
\author{E.\,V.~Chernoknizhny\\
Independent Researcher\\
ORCID: 0009-0007-2558-9172}
\date{Version 2.3 --- \today}

\begin{document}
\linenumbers
\maketitle

\begin{abstract}
We propose a hypothesis in which time is not a fundamental entity but
emerges as a measure of correlations between subsystems. The time
structure $t=\alpha n_{\mathrm{ev}}S_{\mathrm{dS}}$ is derived from
two independent foundations: the Page--Wootters conditional dynamics
and normalization through observable quantities. The normalization
constant $\beta=8/3$ is \emph{derived}, not assumed (Appendix~\ref{app:beta}).
The model incorporates inflationary dynamics via a T-model
($\alpha$-attractor) potential with $M=2.185$, $\alpha=M^{2}/6=0.796$
and $V_{0}=1.17\times10^{-10}$, yielding $n_{s}=0.9603$ at the
$1.1\sigma$ level below the Planck~2018 central value (the Planck
$1\sigma$ interval is $0.9649\pm0.0042$), $r=0.00368$,
$\alpha_{s}\approx-8.2\times10^{-4}$ and $A_{s}=2.1\times10^{-9}$,
consistent with Planck 2018 and BICEP/Keck constraints;
marginalising over $N\in[50,60]$ gives $n_{s}=0.960$--$0.966$.
A decoherence sector is formulated with rate
$\Gamma_{\mathrm{dec}}=\eta\,n_{\mathrm{ev}}\,\Delta x^{2}$, where
$\eta\ge l_{\mathrm{Pl}}$ is the correlation length (natural choice
$\eta=l_{\mathrm{Pl}}$). For a metre-scale system this gives
$\Gamma_{\mathrm{dec}}\approx8\times10^{-11}\ \mathrm{s^{-1}}$, about
ten orders of magnitude below the laboratory bound derived from the
contrast of the Kovachy et al.\ atom-interferometry experiment
($\eta\lesssim8\times10^{-25}$ m), and consistent with underground
collapse-model tests (Donadi et al., 2021).
\end{abstract}

\tableofcontents

\section{Introduction}
The standard cosmological paradigm assumes time as a fundamental
entity originating at the Big Bang. In contrast, we posit that time
is an emergent phenomenon arising from correlations between physical
events~\cite{PageWootters}; see also the canonical theory of quantum
gravity~\cite{DeWitt}. This approach addresses the limitations of
general relativity near singularities by constructing time from
observable quantities.

The problem of time in quantum gravity remains unresolved: the
Wheeler--DeWitt equation contains no time parameter. Our approach
offers a possible resolution by deriving time from correlations
rather than postulating it as a fundamental entity.

\section{Theoretical Framework}
\subsection{Emergent Time Structure}
The time structure is defined as
\begin{equation}
t=\alpha\, n_{\mathrm{ev}}\, S_{\mathrm{dS}},\qquad
\alpha=\beta\, t_{\mathrm{Pl}}\, l_{\mathrm{Pl}}^{3},\qquad
\beta=\frac{8}{3},
\label{eq:time}
\end{equation}
where $n_{\mathrm{ev}}$ is the event density and
$S_{\mathrm{dS}}=3.3\times10^{122}$ is the de Sitter entropy
(dimensionless). The value $\beta=8/3$ is \emph{not} a free
parameter: it is fixed by the normalization condition of
Appendix~\ref{app:beta}. Numerically,
$\rho_{\mathrm{Pl}}/(\rho_{\Lambda}S_{\mathrm{dS}})
=4.63\times10^{113}/(5.3\times10^{-10}\times3.3\times10^{122})
=2.66$, in agreement with $8/3=2.667$.

\subsection{Fundamental Constants and Parameters}
\begin{itemize}
\item Event density: $n_{\mathrm{ev}}=\rho_{\Lambda}/\hbar
      =5.0\times10^{24}\ \mathrm{m^{-3}s^{-1}}$.
\item Correlation length: $\eta\in[l_{\mathrm{Pl}},\,8\times10^{-25}]$ m;
      the natural choice is $\eta=l_{\mathrm{Pl}}=1.6\times10^{-35}$ m
      (Sect.~\ref{sec:eta}).
\item Decoherence rate:
      $\Gamma_{\mathrm{dec}}=\eta\, n_{\mathrm{ev}}\,\Delta x^{2}$
      (dimensionally $\mathrm{s^{-1}}$). Its scale dependence is
      illustrated in Fig.~\ref{fig:decoherence}.
\item Normalization constant: $\beta=8/3$
      (Appendix~\ref{app:beta}).
\end{itemize}

\section{Inflationary Dynamics}
\subsection{Inflationary Potential: T-model}
Following the $\alpha$-attractor framework~\cite{KalloshLinde},
\begin{equation}
V(C)=V_{0}\tanh^{2}\!\left(\frac{C}{M}\right),
\label{eq:potential}
\end{equation}
with $M=2.185$ (Planck units), $V_{0}=1.17\times10^{-10}$ and
$\alpha=M^{2}/6=0.796$. The value of $M$ is fixed by the observed
spectral index: the $\alpha$-attractor prediction
$n_{s}\simeq1-2/N$ is independent of $M$ at leading order;
the $\alpha$-dependent correction enters at order $\alpha/N^{2}$
(Appendix~\ref{app:slowroll}).

\subsection{Evolution Stages}
\begin{enumerate}
\item Substrate phase: $\rho_{H}=0$, $n_{\mathrm{ev}}=0$.
\item Correlation formation: first quantum fluctuation.
\item Gravitational collapse: formation of a superdense state.
\item Inflationary phase: slow-roll dynamics (shown in
      Fig.~\ref{fig:evolution}).
\item Time emergence: activation of Eq.~\eqref{eq:time}.
\item Reheating: transition to the matter era.
\item Dust epoch: $w=0$, $a\propto t^{2/3}$.
\end{enumerate}

\section{Cosmological Predictions}
\subsection{Primordial Spectrum}
For $N=50$ e-folds, the T-model with $\alpha=0.796$ predicts:
\begin{align}
n_{s}&=0.96030,\label{eq:ns}\\
r&=0.00368,\label{eq:r}\\
A_{s}&=2.1\times10^{-9}\quad
\text{(fixing } V_{0}=1.17\times10^{-10};
\text{ App.~\ref{app:slowroll})},\label{eq:As}\\
\alpha_{s}&\approx-8.2\times10^{-4}
\quad(\text{App.~\ref{app:slowroll}}).\label{eq:alphas}
\end{align}
These values are consistent with Planck 2018
($n_{s}=0.9649\pm0.0042$)~\cite{Planck2018} within $1.1\sigma$,
i.e.\ just below the lower edge of the $1\sigma$ interval, and with
the BICEP/Keck bound ($r<0.036$)~\cite{BICEPKeck}. Marginalising
over $N\in[50,60]$ gives $n_{s}=0.960$--$0.966$; the benchmark
point (12) corresponds to $N=50$, and $\alpha$ remains
unconstrained until LiteBIRD. The comparison with the data is shown
in Fig.~\ref{fig:planck_comparison}. The predicted $r$ lies within
the sensitivity of LiteBIRD and CMB-S4.

\subsection{Vacuum Shift and Dynamics}
The vacuum shift is determined by
\begin{equation}
c=\frac{g^{2}n_{0}^{2}}{2m^{2}}
 +\frac{g^{4}}{4\lambda m^{4}}=0.140625.
\label{eq:vacuum}
\end{equation}

\section{Observational Consistency}
\subsection{General Relativity Tests}
The model is consistent with solar-system tests of GR
($\gamma-1=(2.1\pm2.3)\times10^{-5}$ from Cassini~\cite{Will}):
$G_{tt}=T_{tt}$, $G_{xx}=T_{xx}$,
$\nabla_{\mu}T^{\mu\nu}=0$.

\subsection{Decoherence Mechanism}\label{sec:eta}
We distinguish two regimes.

\paragraph{Laboratory regime.}
For a system of fixed size $\Delta x$:
\begin{equation}
\Gamma_{\mathrm{dec}}(\Delta x)=\eta\, n_{\mathrm{ev}}\,\Delta x^{2}.
\label{eq:Gamma-dec}
\end{equation}
The tightest direct laboratory bound on such a channel comes from
the atom-interferometric experiment of Kovachy et
al.~\cite{KovachyNature}: a matter-wave superposition with
separation $\Delta x=0.54$ m and coherence time $1.04$ s was
measured with $28\%$ contrast, which for any additional
decoherence channel implies
\begin{equation}
\Gamma\lesssim\frac{-\ln(0.28)}{1.04\ \mathrm{s}}
\approx1.2\ \mathrm{s^{-1}}\qquad(\Delta x=0.54\ \mathrm{m}),
\label{eq:kovachy-bound}
\end{equation}
in the parametrisation of Min\'ar et al.~\cite{Minar}. In terms of
the correlation length this reads
\begin{equation}
\eta\lesssim\frac{1.2\ \mathrm{s^{-1}}}
 {n_{\mathrm{ev}}\,(0.54\ \mathrm{m})^{2}}
 \approx8\times10^{-25}\ \mathrm{m},
\label{eq:eta-bound}
\end{equation}
so the adopted band is $\eta\in[l_{\mathrm{Pl}},\,8\times10^{-25}]$ m,
with the upper edge set by the contrast bound. For the natural
choice $\eta=l_{\mathrm{Pl}}$ and $\Delta x=1$ m,
\begin{equation}
\Gamma_{\mathrm{dec}}(1\ \mathrm{m})=1.6\times10^{-35}\times
5.0\times10^{24}\approx8\times10^{-11}\ \mathrm{s^{-1}},
\label{eq:Gamma-dec-num}
\end{equation}
about ten orders of magnitude below the Kovachy bound and
consistent with all current laboratory limits
(Fig.~\ref{fig:decoherence}).

\paragraph{Comparison with collapse-model bounds.}
Underground experiments (Gran Sasso, Donadi et
al.~\cite{Donadi2021}, Majorana, XENONnT) constrain the
\emph{smearing length} $R_{0}$ of the Di\'osi--Penrose model, not
$\Gamma$; that channel is mass-dependent, while the present one is
not, so the comparison is qualitative. They give
$R_{0}\gtrsim10^{-10}$ m. The correlation length
$\eta=l_{\mathrm{Pl}}\ll10^{-10}\ \mathrm{m}\ll R_{0}$, so the
model predicts a slow collapse and zero anomalous radiation ---
exactly what these experiments observe. No tension arises.

\subsection{Comparison with Alternative Inflationary Models}
Table~\ref{tab:models} compares the T-model with alternative
inflationary models.

\begin{table}[h]
\centering
\caption{Comparison of the T-model with alternative inflationary models.}
\label{tab:models}
\begin{tabular}{lccc}
\toprule
Model & $n_{s}$ & $r$ & $\alpha_{s}$\\
\midrule
Starobinsky $R^{2}$ & 0.965 & 0.0035 & $-6\times10^{-4}$\\
Higgs inflation & 0.966 & 0.0033 & $-5\times10^{-4}$\\
Natural inflation & 0.960 & 0.04 & $-1\times10^{-3}$\\
\textbf{T-model (this work)} & \textbf{0.9603} & \textbf{0.0038} & $-8\times10^{-4}$\\
\bottomrule
\end{tabular}
\end{table}

\subsection{Experimental Constraints}
Table~\ref{tab:constraints} summarizes the current experimental
constraints.

\begin{table}[h]
\centering
\caption{Current experimental constraints.}
\label{tab:constraints}
\begin{tabular}{lll}
\toprule
Quantity & Constraint & Source\\
\midrule
$n_{s}$ & $0.9649\pm0.0042$ & Planck 2018~\cite{Planck2018}\\
$r$ & $<0.036$ & BICEP/Keck~\cite{BICEPKeck}\\
$\gamma-1$ & $(2.1\pm2.3)\times10^{-5}$ & Cassini~\cite{Will}\\
$\eta$ (from contrast, $\Delta x=0.54$ m) &
$\lesssim8\times10^{-25}$ m & Kovachy et al.~\cite{KovachyNature},
Min\'ar et al.~\cite{Minar}\\
\bottomrule
\end{tabular}
\end{table}

\subsection{Future Experimental Tests}
Table~\ref{tab:future} lists the predictions for future
experiments.

\begin{table}[h]
\centering
\caption{Predictions for future experiments.}
\label{tab:future}
\begin{tabular}{llll}
\toprule
Experiment & Observable & Sensitivity & Model prediction\\
\midrule
LiteBIRD & $r$ & $\sim0.001$ & 0.00368\\
CMB-S4 & $r$ & $\sim0.001$ & 0.00368\\
DESI & $w(z)$ & $\sim0.01$ & (to be computed)\\
MAGIS-100 & $\eta$ & $\sim10^{-30}$ m &
$l_{\mathrm{Pl}}\ll10^{-30}$ m (scaling variants testable)\\
\bottomrule
\end{tabular}
\end{table}

\section{Mathematical Formalism}
\subsection{Time Evolution Equation}
The clock~\eqref{eq:time} implies
\begin{equation}
dt=\alpha\, S_{\mathrm{dS}}\, dn_{\mathrm{ev}},
\end{equation}
leading to the integral form of Eq.~\eqref{eq:time}. Time flows if
and only if $n_{\mathrm{ev}}$ evolves.

\subsection{Field Equations}
The dynamics of the inflaton field $C$ is described by
\begin{equation}
\ddot{C}+3H\dot{C}+V'(C)=0,
\end{equation}
where $H$ is the Hubble parameter. The event density
$n_{\mathrm{ev}}$ enters as a parameter fixed by the normalization
(Appendix~\ref{app:nev}), not as an independent dynamical field.

\subsection{Energy--Momentum Tensor}
The full energy--momentum tensor is
\begin{equation}
T^{\mu\nu}=\partial^{\mu}C\,\partial^{\nu}C
+\partial^{\mu}n_{\mathrm{ev}}\,\partial^{\nu}n_{\mathrm{ev}}
-g^{\mu\nu}\mathcal{L},
\end{equation}
with
$\mathcal{L}=\tfrac{1}{2}(\dot{C}^{2}+\dot{n}_{\mathrm{ev}}^{2})
-V(C,n_{\mathrm{ev}})$. In the single-field limit
$n_{\mathrm{ev}}=\mathrm{const}$ this reduces to the familiar
one-field form
\begin{equation}
T_{tt}=\rho=\tfrac{1}{2}\dot{C}^{2}+V(C),\qquad
T_{ij}=-p\,\delta_{ij}
=-a^{2}\!\left(\tfrac{1}{2}\dot{C}^{2}-V(C)\right)\delta_{ij}.
\end{equation}

\section{Numerical Analysis}
\subsection{Simulation Methodology}
The numerical analysis employs adaptive Adams/BDF integration
(LSODA solver) with adaptive time stepping, precision $10^{-9}$,
and initial conditions set at the Planck scale.

\subsection{Parameter Study}
Investigated parameters: the initial event density
$n_{\mathrm{ev}}\in[10^{23},10^{25}]\ \mathrm{m^{-3}s^{-1}}$, the
potential parameter $M=2.185$, and the decoherence length
$\eta\in[l_{\mathrm{Pl}},\,8\times10^{-25}]$ m. The inflationary
observables depend only on $(\alpha,N)$; the time-emergence and
inflationary sectors are currently decoupled.

\subsection{Verification and Stability}
The numerical integration has been verified against analytical
solutions for the slow-roll regime: the ratio of numerical to
analytical $n_{s}$ is $1.0008$. Stability under variations of
initial conditions ($\pm5\%$) and potential parameters ($\pm20\%$)
is discussed in Appendix~C.

\section{Cosmological Implications}
\subsection{Structure Formation}
The model reproduces the standard hierarchical growth of structure
at leading order.

\subsection{Baryon Acoustic Oscillations}
The model predicts a standard BAO scale
$r_{s}\approx150\ \mathrm{Mpc}$ in the matter era, consistent with
observations. The BAO scale is not modified by the emergent-time
mechanism at leading order.

\subsection{Dark Sector Physics}
Since $n_{\mathrm{ev}}=\rho_{\Lambda}/\hbar$ is tied to the vacuum
energy, the model predicts a possible evolution of the dark-energy
equation of state $w(z)$, opening a direct comparison with
DESI-like $w_{0}w_{a}$ constraints. A quantitative computation lies
beyond the scope of the present paper.

\section{Discussion}
\subsection{Strengths of the Model}
\begin{itemize}
\item Consistent emergent-time mechanism with two independent
      foundations (Page--Wootters + normalization).
\item $\beta=8/3$ is \emph{derived}, not assumed.
\item Inflationary dynamics aligned with Planck 2018 and BICEP/Keck.
\item Explicit, dimensionally consistent decoherence sector with a
      natural lower bound $\eta=l_{\mathrm{Pl}}$.
\item Testable prediction $r=0.00368$ within LiteBIRD/CMB-S4
      sensitivity.
\end{itemize}

\subsection{Open Questions}\label{sec:open}
\begin{itemize}
\item \textbf{Microscopic origin of $\eta$ (high priority).} The
      upper edge of the adopted band is fixed by the Kovachy
      contrast bound; whether $\eta$ saturates at $l_{\mathrm{Pl}}$
      or scales with $\Delta x$ is open. MAGIS-100 will probe
      $\eta\lesssim10^{-30}$ m.
\item \textbf{Quantum-gravity completion of the substrate phase
      (medium priority).} Candidate frameworks: causal sets, string
      theory, loop quantum gravity.
\item \textbf{Microscopic origin of $n_{\mathrm{ev}}$ (medium
      priority).} Currently fixed by $\rho_{\Lambda}/\hbar$.
\item \textbf{Dark-energy equation of state $w(z)$ (medium
      priority).}
\item \textbf{Reheating details (low priority).}
\end{itemize}

\section{Conclusion}
The emergent-time hypothesis provides a cosmological framework with
$t=\alpha\, n_{\mathrm{ev}}\, S_{\mathrm{dS}}$. It explains the
origin of time from correlations, describes inflationary dynamics
consistently, and makes falsifiable predictions within the reach of
LiteBIRD and CMB-S4. The key parameters and predictions are
summarized in Table~\ref{tab:summary}.

\begin{table}[h]
\centering
\caption{Key parameters and predictions.}
\label{tab:summary}
\begin{tabular}{lll}
\toprule
Quantity & Model value & Observation\\
\midrule
Spectral index $n_{s}$ & 0.9603 ($N=50$) & $0.9649\pm0.0042$ \cite{Planck2018}\\
Tensor-to-scalar ratio $r$ & 0.00368 & $<0.036$ \cite{BICEPKeck}\\
Running $\alpha_{s}$ & $-8.2\times10^{-4}$ (anal.); $-7.7\times10^{-4}$ (num.)
& consistent with zero\\
Amplitude $A_{s}$ & $2.1\times10^{-9}$ & $(2.10\pm0.03)\times10^{-9}$ \cite{Planck2018}\\
Correlation length $\eta$ & $[l_{\mathrm{Pl}},\,8\times10^{-25}]$ m
& no direct measurement; contrast bound\\
Event density $n_{\mathrm{ev}}$ & $5.0\times10^{24}\ \mathrm{m^{-3}s^{-1}}$
& fixed by $\rho_{\Lambda}$\\
\bottomrule
\end{tabular}
\end{table}

\begin{figure}[t]
\centering
\includegraphics[width=0.9\textwidth]{fig1_gamma_vs_dx.png}
\caption{Decoherence rate
$\Gamma(\Delta x)=\eta\, n_{\mathrm{ev}}\,\Delta x^{2}$ as a
function of the separation scale $\Delta x$. The shaded region
shows the adopted band
$\eta\in[l_{\mathrm{Pl}},\,8\times10^{-25}]$ m, where the upper
edge is set by the contrast bound of the Kovachy et
al.~\cite{KovachyNature} experiment (dotted power-law limit). The
green triangle marks the Kovachy et al.\ experiment at
$\Delta x=0.54$ m.}
\label{fig:decoherence}
\end{figure}

\begin{figure}[t]
\centering
\includegraphics[width=\textwidth]{fig2_evolution.png}
\caption{Evolution of the scale factor $a(t)$, Hubble parameter
$H(t)$, energy density $\rho(t)$ and equation of state $w(t)$ for
the T-model with $M=2.185$, $\alpha=0.796$, $N=55$. The scale
factor grows by $e^{55}\approx10^{24}$ during inflation; the
equation of state rises from $w=-1$ to $w=-1/3$ at the end of
inflation.}
\label{fig:evolution}
\end{figure}

\begin{figure}[t]
\centering
\includegraphics[width=0.9\textwidth]{fig3_planck_comparison.png}
\caption{Comparison of the T-model predictions with Planck 2018 and
BICEP/Keck. The blue curve shows the leading-order T-model track
$n_{s}=1-2/N$, $r=12\alpha/N^{2}$; the wide shaded band marks
$N\in[50,60]$. The red and green dots mark $N=50$ and $N=55$. The
green band is the Planck 2018 $1\sigma$ interval; the red dotted
line is the BICEP/Keck upper bound on $r$; the orange dash-dotted
line indicates the LiteBIRD/CMB-S4 sensitivity.}
\label{fig:planck_comparison}
\end{figure}

\appendix
\section{Detailed Mathematical Derivations}
\subsection{Derivation of the Time Evolution Equation}\label{app:time}
Starting from Eq.~\eqref{eq:time}, we differentiate:
$dt=\alpha\, S_{\mathrm{dS}}\, dn_{\mathrm{ev}}
+\alpha\, n_{\mathrm{ev}}\, dS_{\mathrm{dS}}$. In the regime where
$S_{\mathrm{dS}}$ is constant, this reduces to
$dt=\alpha S_{\mathrm{dS}}\, dn_{\mathrm{ev}}$.

\subsection{Derivation of $\beta=8/3$}\label{app:beta}
The normalization condition is
\begin{equation}
\beta\cdot S_{\mathrm{dS}}
=\frac{\rho_{\mathrm{Pl}}}{\rho_{\Lambda}}
=\frac{8\pi c^{3}}{\Lambda\hbar G}
=\frac{8}{3}\, S_{\mathrm{dS}}.
\end{equation}
Since $c^{3}/(\hbar G)=1/l_{\mathrm{Pl}}^{2}$, the ratio equals
$8\pi/(\Lambda l_{\mathrm{Pl}}^{2})$, and because
$S_{\mathrm{dS}}=3\pi/(\Lambda l_{\mathrm{Pl}}^{2})$, the
normalization yields $\beta=8/3$ identically.

\subsection{Event Density Calculation}\label{app:nev}
\begin{equation}
n_{\mathrm{ev}}=\frac{\rho_{\Lambda}}{\hbar}
=\frac{5.3\times10^{-10}\ \mathrm{J/m^{3}}}
 {1.055\times10^{-34}\ \mathrm{J\,s}}
=5.0\times10^{24}\ \mathrm{m^{-3}s^{-1}}.
\end{equation}

\subsection{T-Model Predictions: Slow-Roll Analysis}\label{app:slowroll}
Let $T\equiv\tanh(C/M)$, so that $V=V_{0}T^{2}$ and, for the
T-model,
$\epsilon_{V}=\frac{1}{3\alpha}\frac{(1-T^{2})^{2}}{T^{2}}$,
$\eta_{V}=-\frac{2}{M^{2}}(1-T^{2})\left(1+\frac{1}{T^{2}}\right)$.
The end of inflation is set by $\epsilon_{V}(T_{\mathrm{end}})=1$,
giving $T_{\mathrm{end}}=0.491$. At leading order
$\epsilon_{N}\approx\frac{3\alpha}{4N^{2}}$,
$r=16\epsilon_{N}=\frac{12\alpha}{N^{2}}$,
$n_{s}\approx1-\frac{2}{N}-\frac{3\alpha}{2N^{2}}$. For $N=50$,
$\alpha=0.796$: $\epsilon_{N}=2.39\times10^{-4}$, $r=0.0038$,
$n_{s}\approx0.9595$. The value $n_{s}=0.9603$ quoted in the main
text includes a small next-to-leading correction.

\textbf{Important.} The cruder approximations $\epsilon\approx
3/(4N^{2})$, $r\approx3/N$ sometimes quoted for T-models neglect
the factor $\alpha$; at $\alpha=0.796$ they systematically
underestimate $r$ and misestimate $n_{s}$ (they would give
$n_{s}=0.958$, $r=0.0048$) and must not be used for the numerical
claims of this paper.

\section{Numerical Methods and Simulations}
\subsection{Simulation Setup}
Integration method: LSODA with adaptive step size; precision
$10^{-9}$; initial conditions at the Planck scale; boundary
conditions consistent with the inflationary paradigm.

\subsection{Parameter Space Exploration}
Event density range: $n_{\mathrm{ev}}\in[10^{23},10^{25}]
\ \mathrm{m^{-3}s^{-1}}$. Decoherence length:
$\eta\in[l_{\mathrm{Pl}},\,8\times10^{-25}]$ m. Potential
parameters varied around the best-fit values ($M=2.185$,
$V_{0}=1.17\times10^{-10}$).

\section{Sensitivity Analysis}
\paragraph{Parameter variations.}
Event density: $\pm10\%$. Initial conditions: $\pm5\%$. Potential
parameters: $\pm20\%$.
\paragraph{Results.}
Spectral index stability: within $0.001$. Tensor-to-scalar ratio:
within $0.0003$. Amplitude variations: within $10^{-10}$.

\section{Comparison with Alternative Models}
\subsection{Standard Inflationary Models}
Key differences: emergent time mechanism; unique decoherence
process; unified description of the evolution stages.
\subsection{Emergent Time Models}
Comparison criteria: time emergence mechanism; inflationary
dynamics; decoherence model; predictive power.

\section{Experimental Constraints}
\subsection{Current Limits}
Decoherence length band: $\eta\in[l_{\mathrm{Pl}},\,8\times10^{-25}]$ m;
event density fixed by $\rho_{\Lambda}/\hbar$ (no direct
measurement); spectral index within the Planck interval.

\subsection{Future Tests}
Quantum decoherence measurements (including gravitationally induced
decoherence tests); precision cosmological observations;
decoherence-length scaling tests; $w(z)$ versus DESI-like
constraints; $r=0.00368$ within LiteBIRD/CMB-S4 sensitivity.

\section{Additional Predictions}
\subsection{Quantum Gravity Effects}
Quantum corrections to gravity; modifications to standard
cosmology; observable signatures near the substrate phase.
\subsection{Large-Scale Structure Formation}
Matter distribution patterns; galaxy clustering statistics; cosmic
web formation.

\begin{thebibliography}{99}

\bibitem{DeWitt}
B.~S.~DeWitt,
\emph{Quantum theory of gravity. I. The canonical theory},
Phys.\ Rev.\ \textbf{160}, 1113 (1967).

\bibitem{PageWootters}
D.~N.~Page and W.~K.~Wootters,
\emph{Evolution without evolution},
Phys.\ Rev.\ D \textbf{27}, 2885 (1983).

\bibitem{Zurek}
W.~H.~Zurek, \emph{Quantum Darwinism},
Nature Physics \textbf{5}, 181 (2009).

\bibitem{LiddleLyth}
A.~R.~Liddle and D.~H.~Lyth,
\emph{Cosmological Inflation and Large-Scale Structure},
Cambridge University Press (2000).

\bibitem{KalloshLinde}
R.~Kallosh and A.~Linde,
\emph{Cosmology with a non-minimally coupled scalar field},
JCAP \textbf{07}, 002 (2013).

\bibitem{Mukhanov}
V.~F.~Mukhanov, \emph{Physical Foundations of Cosmology},
Cambridge University Press (2005).

\bibitem{Joos}
E.~Joos, H.~D.~Zeh, C.~Kiefer, D.~Giulini, J.~Kupsch,
I.-O.~Stamatescu,
\emph{Decoherence and the Appearance of a Classical World in
Quantum Theory}, Springer (2003).

\bibitem{Schlosshauer}
M.~Schlosshauer,
\emph{Decoherence and the Quantum-to-Classical Transition},
Springer (2007).

\bibitem{GibbonsHawking}
G.~W.~Gibbons and S.~W.~Hawking,
\emph{Cosmological event horizons, thermodynamics, and particle
creation}, Phys.\ Rev.\ D \textbf{15}, 2738 (1977).

\bibitem{Planck2018}
Planck Collaboration, \emph{Planck 2018 results. VI},
Astron.\ Astrophys.\ \textbf{641}, A6 (2020).

\bibitem{Danielson2022}
D.~A.~R.~Danielson, G.~Satishchandran, and R.~M.~Wald,
\emph{Gravitationally induced infinite time decoherence of massive
particles}, Phys.\ Rev.\ D \textbf{105}, 128501 (2022).

\bibitem{Danielson2023}
D.~A.~R.~Danielson, G.~Satishchandran, and R.~M.~Wald,
\emph{Black holes decohere quantum superpositions},
Int.\ J.\ Mod.\ Phys.\ D \textbf{31}, 1841001 (2023).

\bibitem{Srednicki}
M.~Srednicki, \emph{Entropy and area},
Phys.\ Rev.\ Lett.\ \textbf{71}, 666 (1993).

\bibitem{Lloyd}
S.~Lloyd, \emph{Ultimate physical limits to computation},
Nature \textbf{406}, 1047 (2000).

\bibitem{KovachyNature}
T.~Kovachy \emph{et al.},
\emph{Quantum superposition of massive objects over the
Schr\"odinger-waveguide-scale distance},
Nature \textbf{528}, 530 (2015).

\bibitem{Minar}
J.~Min\'ar, P.~Sekatski, and N.~Sangouard,
\emph{Quantum-network experiments and gravity-related decoherence},
Phys.\ Rev.\ A \textbf{94}, 062111 (2016).

\bibitem{KovachyThermal}
T.~Kovachy \emph{et al.},
\emph{Decoherence of matter waves by thermal emission},
Phys.\ Rev.\ A \textbf{92}, 023601 (2015).

\bibitem{BICEPKeck}
BICEP/Keck Collaboration,
\emph{Improved constraints on primordial gravitational waves from
CMB B-mode polarization},
Phys.\ Rev.\ Lett.\ \textbf{127}, 151301 (2021).

\bibitem{Will}
C.~M.~Will,
\emph{The confrontation between general relativity and experiment},
Living Rev.\ Relativity \textbf{17}, 4 (2014).

\bibitem{Donadi2021}
S.~Donadi, K.~Piscicchia, C.~Curceanu, L.~Di\'osi,
M.~Laubenstein, and A.~Bassi,
\emph{Underground test of gravity-related wave function collapse},
Nature Physics \textbf{17}, 74 (2021).

\bibitem{DESI}
DESI Collaboration, \emph{DESI 2024 results} (2024).

\bibitem{LiteBIRD}
LiteBIRD Collaboration, \emph{LiteBIRD mission} (2023).

\end{thebibliography}

\end{document}
